 \setlength{\baselineskip}{20pt}
 \begin{conclusion}

    The thesis focuses on exploring an easy-control, convenient-handling, low-cost, environmentally friendly method to fabricate superhydrophobic surface on the aluminum substrate. On the basis of predecessors’ researches and experiments, the thesis takes Wenzel’s and Cassie-Baxter’s model as the theoretical basis, and takes use of the method of electrochemical etching. The main work of the thesis is showed below:

    (1)	Less pollution to human and environment

    The chemical reagents and chemicals used in the experiments cause no harm to the environment. NaCl is selected as the electrolyte in place of strong acid and strong base. The experiment products Al(OH)3 and H2 are also environmentally friendly. Especially the Al(OH)3 can be used as the reagent in medical treatment. During the modification process, stearic acid is selected as the low surface energy materials in place of fluorosilane, which will cause stimulation to the human skin, eyes and respiratory system.

    (2)	Good controllability of wettability

    The use of electrochemical etching method brings the advantages of multiple control method of wettability. The thesis conducts seven groups of experiment under different conditions, and discussed the effect of current density, electrochemical etching time and concentration of electrolyte on the value of contact angle. And finally the variation tendency of different conditions is given. 

    (3)	Low cost

    The equipments and materials used in this thesis have the advantages of saving the cost. The device is concise, the materials and the reagent are cheap. The electrolyte NaCl can be recycled. These guarantee the low cost of the fabrication. Thus is applicable to mass production.

    However, there is still deficiencies in the thesis, which can be improved in the future.

    (1)	The size limitation

    Due to the capacity of the power source (30V), the effective size of aluminum sheet used in the experiments is 22mm×25mm. If larger area needs to be fabricated, larger power source or longer etching time is a must. 

    (2)	Environmental resistance of superhydrophobicity
    
    The thesis discovers that after some friction or impact the superhydrophobicity will disappear. That means the superhydrophobicity on the aluminum sheet is not that strong. Later research can focus on improving this property.
    

\end{conclusion}


